\section{Conclusion}
\label{sec:conclusion}

We presented \method{}, a training-free retrieval controller whose primary mechanism is marginal-gain early stopping, augmented by conditional adaptive-beam and use-graph levers. The central empirical finding is that the effectiveness of training-free signal-driven adaptation is \emph{judge-dependent}. On a strong 14B judge, \method{} simultaneously improves F1 ($+4.7\%$) and reduces cost ($-15\%$); on a weaker 7B judge the per-split gain is larger (2-hop F1 $+10.3\%$, cost $-23.2\%$, $p{=}0.049$, replicated across three seeds) and a counterfactual force-continue analysis confirms the stopping rule is safe (false-stop rate $7.8\%$). However, on a 7B mixed-depth stream a tuned Fixed-$K{=}2,B{=}8$ \emph{dominates} \method{} and even an Oracle-depth upper bound---an honest boundary showing that when the judge is noisy at depth, a tuned global budget beats per-query adaptation. The cost saving transfers to a second KG (WebQSP/Freebase, $-19\%$) and to the stronger engine. We identify the \emph{judge-quality scope condition} as the key determinant of when training-free GraphRAG controllers succeed, and retract the stronger Pareto-optimality claim an earlier version of this work made. Future work: conformal-calibrated stopping thresholds, frontier-model judges, and natural heterogeneous datasets.
